\section{Mechanization Status}

The Coq artifact currently consists of the following components.

\begin{table}[t]
\caption{Current Coq development status.}
\label{tab:coq-status}
\begin{tabular}{lll}
\toprule
File & Contents & Status \\
\midrule
\texttt{Syntax.v} & source and positioned regex syntax; labeling & implemented \\
\texttt{Sets.v} & finite sets as lists of positions & implemented \\
\texttt{KleeneSemantics.v} & Kleene and marked semantics & partly proved \\
\texttt{PositionAutomaton.v} & nullable/first/last/follow and automata & implemented \\
\texttt{PositionCorrectness.v} & labeling and position lemmas & partly proved \\
\texttt{DegreeofAmbiguity.v} & NFA ambiguity definitions & implemented \\
\texttt{DegreeofInfiniteAmbiguity.v} & Weber--Seidl direction & sketch only \\
\texttt{LindenBridge.v} & bridge to JavaScript regex development & scaffold \\
\bottomrule
\end{tabular}
\end{table}

The most important completed theorem is that the labeling pass preserves the
regular-language semantics after erasing positions.  The current automaton file
contains a marked acceptance predicate for the Glushkov construction, but the
full equivalence between marked regex semantics and marked automaton acceptance
is still open.  The ambiguity file defines accepting-run counts, \(k\)-ambiguity,
finite ambiguity, polynomial ambiguity, and a proof that deterministic automata
viewed as NFAs are unambiguous.

This division gives a clean proof roadmap:
\begin{enumerate}
\item complete marked semantics \(\Leftrightarrow\) position-automaton
acceptance;
\item define the regular-core backtracking tree in Coq;
\item prove Theorem~\ref{thm:trace-run};
\item formalize Weber--Seidl graph patterns as checkable certificates;
\item connect those certificates to ReDoS-oriented warnings.
\end{enumerate}
